\section{Zadanie 17}
n osób \((n \geq 4)\) usiadło formując okrąg. Każda osoba w sposób losowy
wybrała sąsiada. Jaka jest oczekiwana liczba osób niewybranych przez nikogo? A jaka jest
odpowiedź jeśli te osoby siedzą w linii (czyli skrajne osoby mają tylko jednego sąsiada)?

\subsection*{Rozwiązanie}
Najpierw pierwsza część. Niech X - zmienna losowa oznaczająca ilość osób niewybranych przez nikogo. Ponumerujmy osoby na okręgu numerami od 1 do n. Zauważmy, że możemy zapisać:
\[
	X = \sum\limits_{i=1}^{n} X_i,
\]
gdzie \(X_i\) przyjmuje wartość 0, jeśli osoba o numerze i została wybrana przez jednego ze swoich sąsiadów, a wartość 1 w przypadku gdy przez żadnego nie została wybrana. Z liniowości wartości oczekiwanej mamy, że:
\[
	E[X] = E[\sum\limits_{i=1}^{n} X_i] = \sum\limits_{i=1}^{n} E[X_i] = n \cdot \frac{1}{4} = \frac{n}{4}.
\]
ponieważ prawdopodobieństwo że dana osoba nie zostanie wybrana wynosi \(\frac{1}{4}\) - obaj sąsiedzi musieliby wybrać swoich drugich sąsiadów, a to następuje z prawdopodobieństwem \((\frac{1}{2})^2\).

Druga część analogicznie, z tym że numerując sąsiadów od lewej do prawej zauważmy, że \(E[X_1] = E[X_n] = \frac{1}{2}\) oraz \(E[X_2] = E[X_{n-1}] = 0\)(bo zawsze zostaną wybrani przez osoby skrajne) a poza tym, jak wyżej, dla pozostałych indeksów \(E[X_i] = \frac{1}{4}\). Wobec tego:
\[
	E[X] = E[\sum\limits_{i=3}^{n-2} X_i] + E[X_1] + E[X_n] + E[X_2] + E[X_{n-1}] = \sum\limits_{i=3}^{n-2} E[X_i] + 2E[X_n] + 2E[X_{n-1}]
\]
Czyli:
\[
	E[X] = (n-4) \cdot \frac{1}{4} + 2 \cdot \frac{1}{2} + 2 \cdot 0 = \frac{n}{4}
\]

\subsection*{Odpowiedź}
\[
	\boxed{\text{Gdy siedzą w kółku } \frac{n}{4}, \text{gdy w linii } \frac{n}{4} }
\]
